\chapter{Guidance design}\label{ch:guide}

\section{General method}

Building on the idea of classifier guidance we identify following 

Assumption about using the gradient.
We explicitly assume that in order to produce realistic result we have to tie the perturbations to the models' gradient.
We use the none-gradient tied perturbations as baseline for all our experiments.

The hope is also that the gradient is informative in the sense that:
if a method pushes a not secondary target variable to a certain distribution / average value, we would also expect that if we push this variable to the achieved average value, that the previously targeted variable would achieve the same result.
This we explicitly test in test5.

The base guidance algorithm we use works as follows:
...

\section{Guidance loss}

The loss we use to guide the generative model is defined as follows:
...

One of the main contributions of this thesis is the realization that in the weather domain, we can define a guidance target, which can be achieved precisely.
By precisely I mean that the loss can be driven to zero.
While in the image domain, there is always a ground truth during training, at inference we play with different guidance strengths to identify a range where the generated images are good looking.
In the weather domain however, we can define a target that can be closed by tuning the guidance strength.

This is great because instead of tuning the guidance strength to find a range of plausible results, we can define what plausible results look like a priori (guidance target sampling trajectory defined through historical data and plausibility tests), and optimize the guidance strength to achieve this target. 



We essentially have a very similar setting to the guided diffusion for precipitation paper, but we make it more general in these ways:
\begin{itemize}
    \item we define a guidance target explictly, in contrast to observe results across different lambdas
    \item we experiment with different $a_t$ trajectories, in addition to guiding computing the gradient at the first step
    \item we test level of intensity by shifting the intenisty parameter from the flow (lambda) to the trajectory definition ($phi_n$)
    \item Our loss is driven to zero in terms of the difference from the target, instead of being minimized iteratively (a more explicit way to encode the target scenario)
\end{itemize}

\section{Guidance algorithms}

In this section we develop different variants of classifier free guidance, which we adapt to this setting.
We explore the different variants, tune them, expose their flaws, assess their computational requirements, and compare their respective trade off.
In the end we close in on the most promising method in terms of reliability and computational requirements, 
evaluated both in terms of theoretical cost, but also in terms of practical factors such as difficulty of tunining and potential need for multiple runs.

\subsection{CLOSE-GAP}

Greedy

Step is normalized to close the gap as prescribed by the guidance schedule.


\subsection{OPTIMIZE-SCHEDULE}

The step is allowed to take the gradient norm into consideration.
The schedule is optimized to minimize the loss while also complying to a regularization term.
Possible regularization terms include the total guidance.
We can dream up arbitrary penalty terms such as angular defelction, ... .

The problem of this optimization is that the gradient applyed for the guidance at each step changes with respect to the schedule.
The optimization is therefore highly non-convex and very hard to solve.
In fact, the optimizer will converge very fast to a local minimum, highly influenced by the initialization of the schedule.

LBFGS


\subsection{OPTIMIZE-GAIN}

Secant method

The step is allowed to take the gradient norm into consideration.
Since the loss is almost linear with respect to the guidance strength, we can use a secant method to find the guidance strength that minimizes the loss.

In this case we define the secant method as follows:

\subsection{Comparison}

TABLE

Given the analyzed properties of the algorithms, in all our experiments we will make use of the latter presented method.
